\documentclass[10pt]{article}

\usepackage[utf8]{inputenc}

\usepackage[T1]{fontenc}

\usepackage{graphicx}

\usepackage[export]{adjustbox}

\graphicspath{ {./images/} }

\usepackage{amsmath}

\usepackage{amsfonts}

\usepackage{amssymb}

\usepackage[version=4]{mhchem}

\usepackage{stmaryrd}

\usepackage{bbold}

\usepackage{mathrsfs}



\DeclareUnicodeCharacter{0131}{$\imath$}



\begin{document}

Введение в комплексный анализ, 2 изд., ч. 2, М., 1976; [8] II е т-

р е н к о В. П., Рост мероморфных функций, Хар., 1978; [9]



\includegraphics[max width=\textwidth, center]{2023_07_08_0c0299ad06b9518e8489g-1(1)}

Be c k e n ba c E. F., Hut ch is n G.A., "Pacific J. Math. $1969, \mathbf{v}, 28, \mathrm{Na} 1, \mathbf{p} .17-47$.



РАСПРЕДЕЛЕНИЯ ФУНКЦИЯ к а к о Й - ли б о с л у ч a Ï̈ но ї в е ли ч и ны $X$ - функция действительного персменного $x$, принимающая ири каждом $x$ значение, равное вероятности неравенства $X<x$.



Каждая Р. ф. $F(x)$ обладает следующими свойствами:



\begin{enumerate}

  \item $F\left(x^{\prime}\right) \leqslant F\left(x^{\prime \prime}\right)$ при $x^{\prime}<x^{\prime \prime}$;



  \item $F(x)$ непрерывна слева при канкдом $x$;



  \item $\lim F(x)=0, \lim F(x)=1$



\end{enumerate}



\includegraphics[max width=\textwidth, center]{2023_07_08_0c0299ad06b9518e8489g-1(2)}

$X \leqslant x$, и тогда она оказывается непрерывной справа).



В математич. анализе Р. ф. называют любую функцию, для к-рої имеют место своїства 1)-3). Существует взаимно однозначное соответствие между распределе-



\includegraphics[max width=\textwidth, center]{2023_07_08_0c0299ad06b9518e8489g-1}

подмножеств числовой прямої $\mathbb{R}^{1}$ и Р. Ф. Әто соответствие определяется формулой: для любого интервала $[a, b)$



\[

P_{F}([a, b))=F(b)-F(a) .

\]



Каждая функция $F$, обладающая свойствами 1)-3), может рассматриваться как $\mathrm{P}$ ф. нек-рой случайной величины $X$ (напр., случайноїі величины $X(x)=x$, заданной на вероятностном пространстве $\left.\left(\mathbb{R}^{1}, \mathscr{B}, P_{F}\right)\right)$.



Всякая P. ф. может быть однозначно лредставлена в виде суммы



\[

F(x)=a_{1} F_{1}(x)+a_{2} F_{2}(x)+a_{3} b_{3}(x),

\]



где $a_{1}, a_{2}, a_{3}$ - неотрицательные чісла, сумма к-рых равна 1, а $F_{1}, F_{2}, F_{3}-$ Р. ф. такие, что $F_{1}(x)$ абсолютно непрерывна:



\[

F_{1}(x)=\int_{-\infty}^{x} p(z) d z

\]



$F_{2}(x)$ - «ступенчатая функция»:



\[

F_{2}(x)=\sum_{x_{k}<x^{p}}

\]



где $x_{k}$ - точки скачков $F(x)$, а $p_{k}>0$ пропорциональны размеру этих скачков; $F_{3}(x)$ - «сингулярная" компонента - непрерывная функция, производная к-рой почти всюду равна нулю.



П р і м е р. Пусть $X_{k}, k=1,2,3, \ldots$ - бесконечная последовательность независимых случаӥных величін, принимающих значения 1 и 0 с вероятностямі $0<p_{k} \leqslant$ $\leqslant 1 / 2$ и $q_{k}=1-p_{k}$, соответственно. Пусть



\[

X=\sum_{k=1}^{\infty} \frac{X_{k}}{2^{k}},

\]



тогда



\begin{enumerate}

  \item если $p_{k}=q_{k}=1 / 2$ при всех $k$, то $X$ имеет абсолютно непрерывную Р. ф. (с $p(x)=1$ для $0 \leqslant x \leqslant 1$, т. е. $X$ равномерно распределена на $[0,1])$;



  \item если $\sum_{k=1}^{\infty} p_{k}<\infty$, то $X$ имеет «стуненчатую" Р. ф. (она пмеет скачки во всех двоично-рациональных точках отрезка $[0,1])$;



  \item если $\sum_{k=1}^{\infty} p_{k}=\infty, p_{k} \rightarrow 0$ при $k \rightarrow \infty$, то $X$ имеет «сингулярную» Р. ф.



\end{enumerate}



Этот пример является иллюстрацией одної теоремы П. Леві (Р. Lévy), в соответствии с к-рой предел бесконечной свертки дискретных Р. Ф. может содержать только одну из указанных выше комюонент.



«Расстояние» между распределениями $P$ и $Q$ на числовой прямой часто определяют в терминах соответствую- щих Р. .. $F$ и $S$, полагая, напр.,



\[

\rho_{1}(P, Q)=\sup _{x}|F(x)-S(x)|

\]



или



\[

\rho_{2}(P, Q)=\operatorname{Var}(F(x)-S(x))

\]



(см. Распределений сходимость, Леви метрика, Харакперистическая бункция).



Р. ф. наиболее употребительных распределений вероятностей (напр., нормального, биномиального, пуассоновского распределений) табулированы.



Для проверки гипотез о Р. ф. $F$ по результатам независимых наблюдений используют так или иначе измеренное отклонение $F$ от эмпирическоӥ Р. Ф. (см. Колмогорова критерий, Колмогорова - Смирнова критерий, Крамера - Мизеса критерий).



Понятие Р. ф. естественным образом расіространяется на многомерный случай, но многомерные Р. ф. значительно менее употребительны, чем одномерные.



О приближенном представлении Р. ф. см. ГрамаIII арлье ряд, Эджворта ряд, Предельные теоремь



Лum.: [1] К р а м е р Г. Случайные величины и распределения вероятностей, пер. с англ., М., 1947; [2] е г о ж е, Математические методы статистики, пер. с англ., 2 изд., М., 1975; [3] $\Phi$ е л л е р В., Введение в теорию вероятностей и ее приложения, пер. с англ., 2 изд., т. $1-2$, М., 1967; [4] Б о л ь п е в ЈI Н., С м и р н о в Н. В., Таблицы математической статистики,

2 изд., М., 1968.

Ю. В. Прохоров.



РАССЕИВАНИЕ ВЫБОРКИ - одна из скалярных характеристик разброса выборки на прямой относительно какой-либо конкретной точки, называем ой ц е н т р о р р а с с е и в а н іI я, численно равная сумме квадратов отклонениї значений случайных величин, образующих выборку, от центра рассеивания. Пусть $X_{1}, \ldots, X_{n}$ - независимые случайные величины, подчиняющиеся одному и тому же закону, и пусть точка $x, x \in \mathbb{R}^{1}$, выбрана в качестве центра рассеивания. Тогда величина



\[

S_{n}(x)=\sum_{i=1}^{n}\left(X_{i}-x\right)^{2}

\]



наз. р а с с е и в а пи е м в ы б о р ки $X_{1}, \ldots, X_{n}$ относительно центра рассеивания $x$. Так как для любого $x$



\[

S_{n}(x)=S_{n}(\bar{X})+n(\bar{X}-x)^{2} \geqslant S_{n}(\bar{X}),

\]



где $\bar{X}=\frac{1}{n}\left(X_{1}+\ldots+X_{n}\right)$, то Р. в. будет минимальным, если в качестве центра рассеивания выбрать $\bar{X}$. Малые значения Р. в. говорят о сосредоточенности әлементов выборки около центра рассеивания, и наоборот: большие значения Р. в. говорят о большой разбросанності элементов выборки. Понятие Р. в. естественным образом распространяется на многомерные выборки.



Лum.: [1] у и л к с С., Математическая статистика, пер. с англ., М., 1967. М. С. Никулин.



РАССЕИВАНИЯ ЭЛЛИПСОИД - элЛıпсоид в пространстве реализаций случайного вектора, характеризующий сосредоточенность его распределения вероятностей возле нек-рого заданного вектора в терминах моментов 2-го порядка. Пусть случайный вектор $X$, принимающий значения $x=\left(x_{1}, \ldots, x_{n}\right)^{T}$ в $n$-мерном евклидовом пространстве $\mathbb{R} n$, имеет невырожденную ковариационнуло матрицу $B$. В таком случас для любого фиксированного вектора $a, a \in \mathbb{R}^{n}$, в пространстве реализаций $\mathbb{R}^{n}$ можно определить эллипсоид



\[

(x-a)^{T} B^{-1}(x-a)=n-1, \quad x \in \mathbb{R}^{n},

\]



к-рый наз. эл л и пा с о й д о м р а с с е іл в а ни я расıределения вероятностей случайного вектора $X$ относительно вектора $a$, пли эллипсоидом рассеивания случайного вектора $X$. В частности, если $a=\mathrm{E} X$, то Р. э. является геометрич. характеристикой концентра-





\end{document}